% =====================================================================
%  3. MARCO TEÓRICO Y NORMATIVO
%  Responsable: INTEGRANTE 1
%  Extensión objetivo: 6 a 8 páginas — aprox. 2.600 a 3.000 palabras
%  Fuente principal: Sección D de docs/bitacora_decisiones.md
%
%  ADVERTENCIA: esta sección se llena de definiciones y es fácil que quede
%  como un copiar-pegar de normas. Para evitarlo, cierra CADA subsección
%  con una frase que la conecte con el caso Marlin.
% =====================================================================

\section{Marco teórico y normativo}

% TODO (párrafo de entrada, ~80 palabras): anuncia los tres bloques de esta
%   sección: (a) qué es una auditoría ambiental y en qué se diferencia de
%   figuras vecinas; (b) los estándares internacionales aplicables; y (c) el
%   marco peruano de referencia que servirá para el contraste de la sección 7.

% ---------------------------------------------------------------------
\subsection{La auditoría ambiental: definición, tipos y delimitación conceptual}
% Extensión: ~1,5 páginas / 550 palabras

% TODO: Definir auditoría ambiental como proceso SISTEMÁTICO, DOCUMENTADO
%   e INDEPENDIENTE para obtener evidencia objetiva y evaluarla de manera
%   objetiva frente a criterios de auditoría. Subrayar las tres palabras.
%   Fuente: \cite{iso2018iso19011}

% TODO: Tipología por quién audita:
%   - Primera parte (interna): la propia organización se audita.
%   - Segunda parte (externa de proveedor/cliente): audita una parte interesada.
%   - Tercera parte (certificación): organismo independiente acreditado.
%   PREGUNTA CLAVE PARA EL INFORME: ¿de qué tipo fue la HRIA de Marlin?
%   Fue externa e independiente en su ejecución, pero financiada por la
%   empresa auditada. Anticipa aquí en una frase esa tensión y remite a la
%   sección 7. Fuente: \cite{oncommonground2010hria}

% TODO: Tipología por objeto: de sistema de gestión, de cumplimiento legal,
%   de residuos, energética, sectorial vs. integrada.

% TODO: DELIMITACIÓN CONCEPTUAL (este párrafo suele valer varios puntos):
%   - Auditoría   = verificación de conformidad frente a criterios; no sanciona.
%   - Fiscalización / supervisión = potestad ESTATAL de control y sanción
%     (en Perú, el OEFA). Fuente: \cite{congreso2009ley29325}
%   - EIA = instrumento PREDICTIVO y EX ANTE, previo a la operación.
%     Fuente: \cite{congreso2001ley27446}
%   Recomendación: cierra con una tabla comparativa de tres columnas
%   (auditoría / fiscalización / EIA) por filas: quién lo hace, cuándo,
%   carácter vinculante, consecuencia jurídica.

% --- Tabla comparativa (opcional pero muy recomendable) ---------------
%\begin{table}[H]
%  \centering
%  \caption{Diferencias entre auditoría ambiental, fiscalización y evaluación de impacto ambiental}
%  \label{tab:auditoria-vs-fiscalizacion}
%  \begin{tabularx}{\textwidth}{@{}l Y Y Y@{}}
%    \toprule
%    \textbf{Criterio} & \textbf{Auditoría} & \textbf{Fiscalización} & \textbf{EIA} \\
%    \midrule
%    Quién la realiza &  &  &  \\
%    Momento          &  &  &  \\
%    Carácter         &  &  &  \\
%    Consecuencia     &  &  &  \\
%    \bottomrule
%  \end{tabularx}
%  \footnotesize Nota. Elaboración propia a partir de \textcite{iso2018iso19011} y \textcite{congreso2009ley29325}.
%\end{table}

% ---------------------------------------------------------------------
\subsection{ISO 14001:2015 y el sistema de gestión ambiental}
% Extensión: ~1 página / 380 palabras

% TODO: Norma certificable; Estructura de Alto Nivel (Anexo SL) que la hace
%   compatible con ISO 9001 e ISO 45001. Explicar contexto de la
%   organización, liderazgo, planificación, apoyo, operación, evaluación del
%   desempeño y mejora. Señalar que la norma se ORGANIZA sobre el ciclo PHVA,
%   lo que enlaza directamente con la sección 8 de este informe.
%   Fuente: \cite{iso2015iso14001}

% TODO: Concepto de aspecto ambiental vs. impacto ambiental, con un ejemplo
%   tomado del propio caso (aspecto: uso de cianuro en la lixiviación CIL;
%   impacto potencial: contaminación de agua superficial y subterránea).

% ---------------------------------------------------------------------
\subsection{ISO 19011:2018 y la conducción de la auditoría}
% Extensión: ~1,5 páginas / 550 palabras

% TODO: Directrices para auditar sistemas de gestión. Desarrollar:
%   - Los principios: integridad, presentación imparcial, debido cuidado
%     profesional, confidencialidad, independencia, enfoque basado en la
%     evidencia y enfoque basado en riesgos.
%   - Gestión del programa de auditoría.
%   - FASES: inicio -> revisión documental -> preparación del plan ->
%     actividades in situ (reunión de apertura, recolección de evidencia por
%     entrevista, muestreo y observación, reunión de cierre) -> informe ->
%     seguimiento.
%   - Competencia y evaluación de los auditores.
%   Fuente: \cite{iso2018iso19011}

% TODO: DEFINICIONES QUE SE USARÁN EN LA SECCIÓN 6 (déjalas explícitas aquí,
%   porque de ellas depende la clasificación de la matriz de hallazgos):
%   criterio de auditoría, evidencia objetiva, hallazgo, conformidad,
%   NO CONFORMIDAD MAYOR, NO CONFORMIDAD MENOR, OBSERVACIÓN, conclusión de
%   auditoría, acción correctiva y acción preventiva.

% ---------------------------------------------------------------------
\subsection{Estándares internacionales de desempeño ambiental y social}
% Extensión: ~2 páginas / 750 palabras

\subsubsection{IFC Performance Standards}
% TODO: PS1 (evaluación y gestión de riesgos e impactos), PS4 (salud y
%   seguridad de la comunidad), PS5 (adquisición de tierras y reasentamiento
%   involuntario), PS7 (pueblos indígenas, incluye el consentimiento libre,
%   previo e informado) y PS8 (patrimonio cultural).
%   IMPORTANTE: explica por qué son directamente exigibles en este caso —
%   la IFC financió el proyecto con un préstamo de US$ 45 millones, lo que
%   activó sus políticas de salvaguarda. Fuentes: \cite{ifcsfmarlin}, \cite{cao2006marlin}

\subsubsection{Convenio 169 de la OIT y la consulta previa}
% TODO: Obligación de consulta previa, libre e informada a pueblos indígenas
%   y tribales. Es el criterio incumplido más importante de todo el caso.
%   Fuente: \cite{oit1989convenio169}

\subsubsection{Voluntary Principles on Security and Human Rights}
% TODO: Marco para empresas extractivas sobre uso de seguridad privada y
%   pública. Se recomendó su adopción justamente por el eje de seguridad
%   de la auditoría. Fuente: \cite{vpshr2000principios}

\subsubsection{ICMM, GISTM e IRMA}
% TODO: - ICMM Mining Principles. Fuente: \cite{icmmsfprincipios}
%       - GISTM (2020): desarrollado por ICMM, UNEP y PRI DESPUÉS de
%         Brumadinho; seis áreas temáticas, 15 principios y 77 requisitos
%         auditables; su Tema I exige debida diligencia en derechos humanos
%         y participación de las personas afectadas.
%         Fuentes: \cite{icmm2020gistm}, \cite{unep2020gistm}
%       - IRMA: más de 420 requisitos auditables, auditoría independiente de
%         tercera parte, y exigencia de al menos 90 % de cumplimiento en cada
%         una de las cuatro áreas de principios para certificar.
%         Fuente: \cite{irmasfstandard}

% ---------------------------------------------------------------------
\subsection{Marco normativo peruano de referencia}
% Extensión: ~1,5 páginas / 550 palabras
% NOTA: esta subsección NO se aplica al caso guatemalteco. Sirve para el
% contraste de la sección 7.5. Dilo explícitamente para que no parezca error.

% TODO: - Ley 28611, Ley General del Ambiente. \cite{congreso2005ley28611}
%       - Ley 27446, Ley del SEIA, y su reglamento D.S. 019-2009-MINAM;
%         rol del SENACE en la aprobación de EIA detallados.
%         \cite{congreso2001ley27446}, \cite{minam2009ds019}
%       - Ley 29325, SINEFA: el OEFA como ente rector de la fiscalización,
%         supervisión, evaluación, control y sanción ambiental de la mediana
%         y gran minería. \cite{congreso2009ley29325}
%       - D.S. 040-2014-EM, Reglamento de Protección y Gestión Ambiental
%         Minera; medidas correctivas del artículo 39. \cite{minem2014ds040}
%       - D.S. 042-2003-EM, compromiso previo y relación con comunidades.
%         \cite{minem2003ds042}
%       - Ley 29785, Ley de Consulta Previa. \cite{congreso2011ley29785}
%       - Mención de OSINERGMIN (seguridad de infraestructura), SUNAFIL
%         (seguridad y salud en el trabajo) y de la figura histórica del PAMA.

% ---------------------------------------------------------------------
\subsection{El ciclo PHVA como base de la mejora continua}
% Extensión: ~0,5 página / 200 palabras

% TODO: Planificar - Hacer - Verificar - Actuar, aplicado a la gestión
%   ambiental y social minera. Explicar brevemente cada fase porque la
%   sección 8 se construye enteramente sobre esta estructura:
%   - Planificar: política, aspectos e impactos, objetivos y metas.
%   - Hacer: controles operacionales, competencias, comunicación.
%   - Verificar: monitoreo, medición, auditoría interna, monitoreo
%     participativo (modelos AMAC en Marlin y mesa de diálogo en Espinar).
%   - Actuar: acciones correctivas, revisión por la dirección, cierre de
%     no conformidades.
%   Fuente: \cite{iso2015iso14001}

\pendiente{desarrollar la sección 3 (6--8 páginas)}
